\section*{Chapitre \ref{ch:fm:the-proprietary-market-making-firm} --- L'entreprise de market making pour compte propre}

\begin{solution}{exo:fm:the-proprietary-market-making-firm:1}
Minimum permanent 0.75~million de dollars~; \omterm{def:fm:the-proprietary-market-making-firm:ownfunds}{exigence au titre des frais généraux fixes} $0.25\times95=\$23.75$ millions~; \omterm{def:fm:the-proprietary-market-making-firm:ownfunds}{exigence au titre des facteurs K} (K-DTF) $0.1\%\times10\,000=\$10$ millions (aucune
charge sur la position nette). C'est l'\omterm{def:fm:the-proprietary-market-making-firm:ownfunds}{exigence au titre des frais généraux fixes} qui est contraignante~: 23.75~millions de dollars.
\end{solution}

\begin{solution}{exo:fm:the-proprietary-market-making-firm:2}
\omterm{def:fm:the-proprietary-market-making-firm:ownfunds}{Exigence au titre des facteurs K} (K-DTF) $0.01\%\times10\,000=\$1$ million. L'exigence reste de 23.75~millions de dollars, fixée par les frais généraux fixes.
\end{solution}

\begin{solution}{exo:fm:the-proprietary-market-making-firm:3}
Allocations 46.5, 34.875, 23.25 et 11.625~; prélèvements de 37.2, 27.9, 18.6 et 9.3~millions de dollars, soit 93.0~millions de dollars au total~; 23.25~millions de dollars restent dans l'entreprise.
\end{solution}

\begin{solution}{exo:fm:the-proprietary-market-making-firm:4}
$g+d=35\%$ du stock chaque année (la \cref{prop:fm:the-proprietary-market-making-firm:parity})~; les dépenses croissent de 15~\% par an~: $35\times1.15^9=\$123$
millions la dixième année.
\end{solution}

\begin{solution}{exo:fm:the-proprietary-market-making-firm:5}
Revenu du marché complet $10^{-4}\times20\,000\times250=\$500$ millions. Le bénéfice avant rémunération variable est nul lorsque $35\times1.15^{t-1}=500-60=440$~:
$t-1=\ln(440/35)/\ln1.15=18.1$, soit la vingtième année.
\end{solution}

\begin{solution}{exo:fm:the-proprietary-market-making-firm:6}
Au-delà de 20~\%, l'entreprise sature son marché en quelques années et continue de retenir du capital qu'elle ne peut pas utiliser, ce qui retarde les versements sans aucun gain. Tout taux fixe retient
soit trop peu tant que le capital manque, soit trop une fois qu'il n'en manque plus~; la politique de remplissage puis de versement est optimale parmi les politiques dépendant du capital
(la \cref{prop:fm:the-proprietary-market-making-firm:retain}).
\end{solution}

\begin{solution}{exo:fm:the-proprietary-market-making-firm:7}
La cinquante-sixième année. Les dépenses de parité croissent de 15~\% par an et le revenu, au plus de 10~\% comme le marché, si bien que les dépenses rattrapent à la fin le revenu diminué des coûts fixes,
quelle que soit la marge de départ. Une \omterm{def:fm:the-proprietary-market-making-firm:treadmill}{course technologique} n'est soutenable que tant que le marché de l'entreprise croît au moins aussi vite que les dépenses de ses concurrents.
\end{solution}

\begin{solution}{exo:fm:the-proprietary-market-making-firm:8}
La réduction des dépenses fait économiser immédiatement et ne fait perdre de la capture qu'à mesure que le stock se déprécie par rapport aux concurrents~: dans le modèle, l'entreprise plus sobre gagne davantage
la première année (124 millions de dollars contre 116 millions) et moins à partir de la deuxième, la capture atteignant 0.44 de la parité à la dixième année. Il faut mesurer la capture par unité échangée et
la vitesse relative, et non le bénéfice d'une seule année.
\end{solution}

\begin{solution}{pb:fm:the-proprietary-market-making-firm:1}
\begin{enumerate}
\item Capital réglementaire, pour l'essentiel des \omterm{def:fm:the-proprietary-market-making-firm:ownfunds}{fonds propres} et des bénéfices mis en réserve~; l'exigence est la plus élevée entre le minimum permanent et un quart
des frais généraux fixes de l'année précédente et les facteurs K.
\item $0.25\times95=\$23.75$ millions.
\item 10~millions de dollars sur les transactions au comptant~; 1~million de dollars sur le notionnel des dérivés.
\item L'\omterm{def:fm:the-proprietary-market-making-firm:ownfunds}{exigence au titre des frais généraux fixes}, 23.75~millions de dollars~; les prime brokers en exigent 200~millions de dollars, soit 8.4 fois ce montant.
\item Le capital net, c'est-à-dire la valeur nette diminuée des actifs peu liquides et des décotes sur les positions, doit rester au-dessus d'un plancher fixé par l'activité et par l'endettement ou
les débits des clients.
\item Une personne morale dont les membres se partagent le bénéfice selon un accord et dont la responsabilité est limitée~; le capital qu'un associé a apporté ou laissé dans
la société, remboursable à son départ.
\item 10~milliards de dollars par jour, 250~millions, 35~millions et 116.25~millions de dollars.
\item Allocation de 46.5~millions de dollars~; tirage de 37.2~millions de dollars~; 9.3~millions de dollars restent dans la société.
\item À la fin de la huitième année, contre la sixième année sans le départ.
\item La part du bénéfice conservée sous forme de capital.
\item 582, 939, 706 et 296 millions de dollars.
\item 20~\%, 939~millions de dollars.
\item 1\,210~millions de dollars, soit 271~millions de plus~: le capital rapporte beaucoup tant que le marché n'est pas saturé, et rien ensuite.
\item Les sixième, troisième et deuxième années.
\item Le revenu dépend de la technologie par rapport à celle des concurrents, si bien que l'entreprise doit maintenir ses dépenses au rythme de ces derniers pour simplement rester sur place.
\item $s=(g+d)T$~: 35~\% du stock.
\item 0.44 de la parité.
\item La vingtième année pour l'entreprise qui dépense à parité, la dix-huitième pour celle qui dépense moins.
\item Les dépenses de parité représentent 35~\% du stock technologique, 35~millions de dollars la première année et 123~millions de dollars la dixième~; tout retenir jusqu'à ce que le capital
sature le marché, puis tout verser, vaut 1\,210~millions de dollars contre 939~millions pour le meilleur taux fixe (20~\%).
\item Garder le bénéfice jusqu'à ce que le capital de trading de l'entreprise sature le marché qu'elle peut servir, puis verser~; maintenir les dépenses technologiques au rythme des concurrents
et chercher des marchés qui croissent au moins aussi vite.
\end{enumerate}
\end{solution}

\begin{solution}{iq:fm:the-proprietary-market-making-firm:1}
Son \omterm{def:fm:the-proprietary-market-making-firm:ownfunds}{exigence au titre des frais généraux fixes} correspond à un quart des frais généraux fixes de l'année précédente~: l'embauche et la technologie la font monter quoi que l'entreprise négocie.
\iqlookfor{l'\omterm{def:fm:the-proprietary-market-making-firm:ownfunds}{exigence au titre des frais généraux fixes} et le fait qu'elle peut devenir contraignante avant le risque.}
\end{solution}

\begin{solution}{iq:fm:the-proprietary-market-making-firm:2}
Généralement, celles des prime brokers et des sociétés de compensation~: elles exigent un capital déposé en garantie des positions, selon leurs modèles de décote et de marge, souvent
plusieurs fois le montant réglementaire.
\iqlookfor{les exigences de dépôt fixent la capacité de trading.}
\end{solution}

\begin{solution}{iq:fm:the-proprietary-market-making-firm:3}
Pas la totalité~: tant que le capital reste inférieur à ce que le marché peut absorber, chaque dollar mis en réserve rapporte le rendement du capital de l'entreprise, très supérieur à l'alternative des
associés. Il faut en garder assez pour atteindre 400~millions de dollars, puis verser.
\iqlookfor{le rendement marginal du capital mis en réserve et le plafond de capacité.}
\end{solution}

\begin{solution}{iq:fm:the-proprietary-market-making-firm:4}
$g+d=35~\%$ du stock par an, et le montant croît de 10~\%.
\iqlookfor{la condition de régime permanent $s=(g+d)T$.}
\end{solution}

\begin{solution}{iq:fm:the-proprietary-market-making-firm:5}
Le capital qui sort est du capital de trading~: le rembourser d'un coup réduirait le volume de l'entreprise et pourrait la faire manquer à ses exigences. Échelonner le remboursement transforme le
départ en une réduction planifiée.
\iqlookfor{le capital comme capacité, et le plancher réglementaire.}
\end{solution}

\begin{solution}{iq:fm:the-proprietary-market-making-firm:6}
La capture par unité échangée, la latence relative et la qualité des exécutions sur plusieurs trimestres~: l'économie est immédiate, alors que la perte de capture arrive à mesure que la
technologie prend du retard sur celle des concurrents.
\iqlookfor{le décalage entre les dépenses et la capture, et les mesures relatives.}
\end{solution}
